\chapter{November 2025 paper}
\label{ch:paper-2025-nov}

\begin{whatfor}
The end-of-semester paper of November 2025 (3 hours, 50 marks), solved question by question. Every solution names its card and row, says what the card says in plain words, works one substitution per step, and ends with the answer.
\end{whatfor}

%=====================================================================
\section{Question 1: short answers}

Throughout, $\{W_t,\ t\ge0\}$ is a standard Wiener process. The answer table for the first page:

\begin{center}
\begin{tabular}{@{}cl@{\qquad}cl@{}}
\toprule
Part & Answer & Part & Answer\\
\midrule
(a) & $\frac19\approx0.11$ & (d) & $\frac23\approx0.67$\\
(b) & $e^tW_t\,dt+e^t\,dW_t$ & (e) & Rs $4.39$\\
(c) & No & &\\
\bottomrule
\end{tabular}
\end{center}

\begin{problem}[Q1(a)]
Compute the variance of $\int_0^1t^4\,dW_t$.
\uses{card:itoint}{2}
\cardsays The Itô integral has mean $0$, and its variance is the ordinary integral of the integrand squared (the Itô isometry):
\[
\Var\int_0^tb_s\,dW_s=\int_0^t\E_P(b_s^2)\,ds.
\]
\working
\begin{steps}
\item Integrand: $b_s=s^4$, not random, so $\E_P(b_s^2)=s^8$.
\item Integrate: $\int_0^1s^8\,ds=\Big[\frac{s^9}9\Big]_0^1=\frac19$.
\end{steps}
\answer{$\frac19\approx0.11$.}
\end{problem}

\begin{problem}[Q1(b)]
Apply Itô's formula to compute $d(e^tW_t)$.
\uses{card:ito}{1}
\cardsays For a function of time and $W$: $df(t,W_t)=\big(\partial_tf+\frac12\partial_{xx}f\big)dt+\partial_xf\,dW_t$.
\working
\begin{steps}
\item The function: $f(t,x)=e^tx$.
\item Derivatives: $\partial_tf=e^tx$, $\partial_xf=e^t$, $\partial_{xx}f=0$.
\item Substitute, with $x=W_t$: $dt$-part $e^tW_t+0$; $dW$-part $e^t$.
\end{steps}
\answer{$d(e^tW_t)=e^tW_t\,dt+e^t\,dW_t$.}
\end{problem}

\begin{problem}[Q1(c)]
Is $Y_t=e^tW_t$ a martingale?
\uses{card:ito}{3}
\cardsays $g(t,W_t)$ is a martingale exactly when its $dt$-coefficient $\partial_tg+\frac12\partial_{xx}g$ is $0$ for every $x$.
\working
\begin{steps}
\item From (b): the $dt$-coefficient is $e^tW_t$.
\item It is not $0$ for every value of $W_t$, so the test fails.
\item Check by the mean: $\E_P(e^tW_t\given\F_s)=e^tW_s\ne e^sW_s$ when $t\ne s$.
\end{steps}
\answer{No.}
\end{problem}

\begin{problem}[Q1(d)]
$T=\min\{t:W_t=2-3t\}$, the first time $W$ hits the line $2-3t$. Use the martingale stopping theorem to find $\E_PT$.
\uses{card:ost}{5}
\cardsays A stopped martingale keeps its start, $\E_PM_\tau=M_0$. With $M=W$ (start $0$) and a line $c+\beta t$ moving towards the path, $W_\tau=c+\beta\tau$, so
\[
0=c+\beta\,\E_P\tau\qquad\Longrightarrow\qquad\E_P\tau=-\frac c\beta.
\]
\working
\begin{steps}
\item Read the line $2-3t$: $c=2$, $\beta=-3$.
\item Guard: $c>0$, $\beta<0$, opposite signs, so the line comes down to meet the path.
\item Substitute: $\E_PT=-\frac{2}{-3}=\frac23$.
\end{steps}
\answer{$\E_PT=\frac23\approx0.67$.}
\end{problem}

\begin{problem}[Q1(e)]
A European call: stock Rs $80$, strike Rs $85$, expiry $1$ year, risk-free rate $5\%$, volatility $\sigma=0.15$. Find its value.
\uses{card:bsprice}{1}
\cardsays The Black--Scholes call price, with $\tau$ the time left:
\[
C=S\,\Phi(h)-Ke^{-r\tau}\,\Phi(h-\sigma\sqrt\tau),\qquad h=\frac{\ln(S/K)+(r+\frac12\sigma^2)\tau}{\sigma\sqrt\tau}.
\]
\working
\begin{steps}
\item Read off: $S=80$, $K=85$, $r=0.05$, $\sigma=0.15$, $\tau=1$.
\item Pieces: $\ln\frac{80}{85}=-0.06062$; $r+\frac12\sigma^2=0.06125$; $\sigma\sqrt\tau=0.15$.
\item $h=\frac{-0.06062+0.06125}{0.15}=0.0042$ and $h-\sigma\sqrt\tau=-0.1458$.
\item Table: $\Phi(0.0042)=0.5017$ and $\Phi(-0.1458)=0.4420$.
\item Discounted strike: $85e^{-0.05}=80.85$.
\item Substitute: $C=80(0.5017)-80.85(0.4420)=40.13-35.74=4.39$.
\end{steps}
\answer{The call is worth Rs $4.39$.}
\end{problem}

%=====================================================================
\section{Question 2: Greeks, parity, no arbitrage, straddle, value at risk}

\begin{problem}[Q2(a)]
Interpret the Greeks delta, gamma and theta of a European call in the Black--Scholes framework.
\uses{card:bsprice}{6}
\cardsays Each Greek is a derivative of the call price $C=S\Phi(h)-Ke^{-r\tau}\Phi(h-\sigma\sqrt\tau)$:
\[
\Delta=\Phi(h),\qquad\Gamma=\frac{\varphi(h)}{S\sigma\sqrt\tau},\qquad\theta=-\frac{S\varphi(h)\sigma}{2\sqrt\tau}-rKe^{-r\tau}\Phi(h-\sigma\sqrt\tau).
\]
\working
\begin{steps}
\item \textbf{Delta} $=\frac{\partial C}{\partial S}=\Phi(h)$, between $0$ and $1$. Meaning: if the stock rises by Rs $1$, the call rises by about $\Delta$. It is also the hedge: hold $\Delta$ shares per call sold to cancel small moves. Deep in the money $\Delta\to1$; far out of the money $\Delta\to0$.
\item \textbf{Gamma} $=\frac{\partial^2C}{\partial S^2}=\frac{\varphi(h)}{S\sigma\sqrt\tau}>0$. Meaning: how fast delta changes, so how often the hedge must be adjusted. It is largest near $S=K$ close to expiry. $\Gamma>0$ says the call price is convex in $S$.
\item \textbf{Theta} $=\frac{\partial C}{\partial t}<0$ for a call. Meaning: the value lost per unit of time as expiry comes closer, with everything else fixed (time decay).
\item Numbers from Q1(e) ($S=80$, $K=85$, $h=0.0042$, $\varphi(h)=0.3989$): $\Delta=0.5017$; $\Gamma=\frac{0.3989}{80(0.15)(1)}=0.0332$; $\theta=-\frac{80(0.3989)(0.15)}{2}-0.05(80.85)(0.4420)=-2.39-1.79=-4.18$ per year.
\end{steps}
\answer{Delta: price change per Rs $1$ of stock, and the hedge ratio. Gamma: change in delta per Rs $1$, the curvature. Theta: value lost per unit of time.}
\end{problem}

\begin{problem}[Q2(b)]
Derive put--call parity.
\uses{card:parity}{3}
\cardsays For European options on the same stock with the same $K$ and $T$, rate $r$ compounded continuously:
\[
C-P=S_0-Ke^{-rT}.
\]
\working
\begin{steps}
\item Portfolio A: one call and cash $Ke^{-rT}$ in the bond (it grows to $K$ at $T$). Portfolio B: one put and one share.
\item At $T$, if $S_T>K$: A is worth $(S_T-K)+K=S_T$; B is worth $0+S_T=S_T$.
\item At $T$, if $S_T\le K$: A is worth $0+K=K$; B is worth $(K-S_T)+S_T=K$.
\item So both always pay $\max(S_T,K)$. (In one line: $(S_T-K)^+-(K-S_T)^+=S_T-K$, \cref{card:payoff}.)
\item Same payoff, so the same price today, or there is an arbitrage: $C+Ke^{-rT}=P+S_0$.
\item If it fails, say $C+Ke^{-rT}>P+S_0$: sell A, buy B, and keep the difference; at $T$ the two cancel exactly.
\end{steps}
\answer{$C+Ke^{-rT}=P+S_0$, that is $C-P=S_0-Ke^{-rT}$.}
\end{problem}

\begin{problem}[Q2(c)]
Argue that the Black--Scholes market is arbitrage-free.
\uses{card:bsprice}{7}
\cardsays The words used below.
\begin{itemize}
\item A \emph{trading strategy} $(\Delta_t,b_t)$ holds $\Delta_t$ shares and $b_t$ bonds; its value is $V_t=\Delta_tS_t+b_tB_t$. It is \emph{self-financing} if it changes only by the moves of the prices: $V_{t+1}-V_t=\Delta_{t+1}(S_{t+1}-S_t)+b_{t+1}(B_{t+1}-B_t)$ in discrete time, $dV_t=\Delta_t\,dS_t+b_t\,dB_t$ in continuous time.
\item An \emph{arbitrage} is a self-financing strategy with $V_0=0$, $V_T\ge0$ always, and $P(V_T>0)>0$. The market is \emph{arbitrage-free} if there is none.
\item An \emph{EMM} (equivalent martingale measure) is a probability $P^*$ that gives positive probability to the same events as $P$, under which the discounted stock $\tilde S_t=S_t/B_t$ is a martingale.
\item A claim $X$ is \emph{attainable} if some self-financing strategy has $V_T=X$. The market is \emph{complete} if every claim is attainable.
\end{itemize}
Two key facts: if an EMM exists, the market is arbitrage-free; if the EMM is unique, the market is complete. Below, each of the four statements is proved directly.
\working

\paragraph{The Black--Scholes market.} Bond $B_t=e^{rt}$; stock $dS_t=\mu S_t\,dt+\sigma S_t\,dW_t$ with $\mu\in\mathbb R$, $\sigma>0$, $t\in[0,T]$.
\begin{steps}
\item First, find every EMM. An EMM is a new probability $Q$ for the same paths. So ask: when we switch from $P$ to $Q$, what about the stock can change?
\begin{deriv}
\dl{d\ln S_t=\big(\mu-\tfrac12\sigma^2\big)dt+\sigma\,dW_t}{\cref{card:ito}, row 5, with $f=\ln x$}
\dl{\ln S_t=\ln S_0+\big(\mu-\tfrac12\sigma^2\big)t+\sigma W_t}{integrate from $0$ to $t$}
\dl{\textstyle\sum\big(\Delta\ln S\big)^2\to\sigma^2t}{over finer and finer grids on $[0,t]$; $(dW)^2=dt$ (\cref{card:itoint}): this holds on every path, with $P$-probability $1$}
\dl{\text{the same under }Q}{$Q$ is equivalent: an event of probability $1$ under $P$ keeps probability $1$ under $Q$}
\dl{\sigma\text{ cannot change}}{$\sigma$ can be read off a single path}
\dl{W_t=W^Q_t+\int_0^t\theta_s\,ds}{the drift cannot be read off a path, and it can change: under $Q$, $W$ is a Brownian motion $W^Q$ plus a drift $\theta_t$ (adapted) and nothing else (Girsanov's theorem)}
\dl{\ln S_t=\ln S_0+\int_0^t\big(\mu-\tfrac12\sigma^2+\sigma\theta_s\big)ds+\sigma W^Q_t}{substitute into the second line}
\dl{S_t=S_0\exp\Big\{\int_0^ta_s\,ds+\sigma W^Q_t\Big\}}{name the new drift $a_s=\mu-\frac12\sigma^2+\sigma\theta_s$}
\end{deriv}
So each candidate $Q$ gives the stock the same $\sigma$ and some adapted drift $a_t$. The question ``which $Q$ is an EMM?'' becomes ``which drift $a_t$ makes $e^{-rt}S_t$ a martingale?''. From here on write $W$ for $W^Q$ and $Z_t$ for the stock under $Q$:
\begin{deriv}
\dl{Z_t=S_0\exp\Big\{\int_0^ta_s\,ds+\sigma W_t\Big\}}{the stock under a candidate $Q$; under $P$ itself, $a_t=\mu-\frac12\sigma^2$}
\dl{e^{-rt}Z_t=S_0X_tY_t}{split into two factors}
\dl{X_t=\exp\Big\{\int_0^ta_s\,ds\Big\},\qquad Y_t=e^{-rt+\sigma W_t}}{$X_t$ has no $dW$ part}
\end{deriv}
\item Differentiate each factor:
\begin{deriv}
\dl{dX_t=a_tX_t\,dt}{derivative of $e^{\int_0^ta_s\,ds}$ in $t$}
\dl{dY_t=Y_t\big(-r+\tfrac12\sigma^2\big)dt+\sigma Y_t\,dW_t}{\cardrow{card:ito}{2} with $\mu=-r$}
\dl{dX_t\,dY_t=0}{$dt\cdot dt=dt\cdot dW=0$}
\end{deriv}
\item The product rule:
\begin{deriv}
\dl{d(X_tY_t)=Y_t\,dX_t+X_t\,dY_t+dX_t\,dY_t}{\cardrow{card:ito}{6}}
\dl{\phantom{d(X_tY_t)}=X_tY_t\big(a_t-r+\tfrac12\sigma^2\big)dt+\sigma X_tY_t\,dW_t}{substitute Step~2; collect the $dt$ terms}
\end{deriv}
\item The martingale test:
\begin{deriv}
\dl{e^{-rt}Z_t\text{ is a martingale}}{what an EMM needs}
\dl{\iff X_tY_t\big(a_t-r+\tfrac12\sigma^2\big)=0}{\cref{card:ito}, martingale test: the $dt$-coefficient is $0$}
\dl{\iff a_t=r-\tfrac12\sigma^2\ \text{ for all }t}{$X_tY_t>0$: both are exponentials}
\end{deriv}
So there is \textbf{exactly one} EMM $P^*$: the one under which the stock is $S_0e^{(r-\frac12\sigma^2)t+\sigma W^*_t}$, with $W^*$ a Brownian motion under $P^*$. From Step~3 with this $a_t$:
\begin{deriv}
\dl{d\big(e^{-rt}S_t\big)=\sigma e^{-rt}S_t\,dW^*_t}{the $dt$ term is gone}
\end{deriv}
\end{steps}

\paragraph{Derivation 1: the Black--Scholes market is arbitrage-free.}
\begin{steps}[resume]
\item Any self-financing strategy, discounted. Write $\tilde V_t=e^{-rt}V_t$:
\begin{deriv}
\dl{d\tilde V_t=e^{-rt}\big(dV_t-rV_t\,dt\big)}{product rule; $d(e^{-rt})=-re^{-rt}dt$}
\dl{\phantom{d\tilde V_t}=e^{-rt}\big(\Delta_t\,dS_t+b_trB_t\,dt}{self-financing, $dB_t=rB_t\,dt$,}
\dl{\phantom{d\tilde V_t={}}-r\Delta_tS_t\,dt-rb_tB_t\,dt\big)}{and $V_t=\Delta_tS_t+b_tB_t$}
\dl{\phantom{d\tilde V_t}=\Delta_t\,e^{-rt}\big(dS_t-rS_t\,dt\big)=\Delta_t\,d\big(e^{-rt}S_t\big)}{the $b_t$ terms cancel}
\dl{\phantom{d\tilde V_t}=\Delta_t\,\sigma e^{-rt}S_t\,dW^*_t}{Step~4}
\end{deriv}
\item So $\tilde V_t$ is a $P^*$-martingale:
\begin{deriv}
\dl{\tilde V_T=\tilde V_0+\int_0^T\Delta_t\sigma e^{-rt}S_t\,dW^*_t}{integrate Step~5}
\dl{\E_{P^*}\tilde V_T=\tilde V_0=V_0}{an Itô integral has mean $0$ (\cref{card:itoint}); $\Delta_t$ with $\int_0^T\E_{P^*}(\Delta_tS_t)^2dt<\infty$}
\end{deriv}
\item Suppose there were an arbitrage. Then:
\begin{deriv}
\dl{\E_{P^*}\big(e^{-rT}V_T\big)=V_0=0}{Step~6, with $V_0=0$}
\dl{V_T\ge0\ \text{and}\ P^*(V_T>0)>0}{$P^*$ and $P$ agree on which events have probability $0$}
\dl{\E_{P^*}\big(e^{-rT}V_T\big)>0}{a variable that is $\ge0$ and sometimes $>0$ has positive mean}
\end{deriv}
The last two lines contradict each other. \textbf{So there is no arbitrage.} (This is the same proof as for one period: an EMM gives every self-financing strategy mean discounted gain $0$.)
\end{steps}

\paragraph{Derivation 2: the Black--Scholes market is complete.}
\begin{steps}[resume]
\item By the uniqueness found in Step~4 and the second key fact, the market is complete. Concretely, here is the copy of a claim $X=g(S_T)$:
\begin{deriv}
\dl{P_t(X)=\E_{P^*}\big(e^{-r(T-t)}g(S_T)\given\F_t\big)}{the candidate price at time $t$}
\dl{\phantom{P_t(X)}=e^{-r(T-t)}\E_{P^*}\big(g(S_T)\given S_t\big)=:f(t,S_t)}{the stock is Markov: only today's price matters}
\dl{df(t,S_t)=\big(\partial_tf+\tfrac12\sigma^2S_t^2\partial_{ss}f\big)dt+\partial_sf\,dS_t}{\cref{card:ito}, row 5; $(dS_t)^2=\sigma^2S_t^2\,dt$}
\dl{dV_t=\Delta_t\,dS_t+b_t\,dB_t}{what a self-financing copy must do}
\end{deriv}
\item Match the two lines:
\begin{deriv}
\dl{\Delta_t=\partial_sf(t,S_t)}{the $dS_t$ terms}
\dl{b_t=\frac{f(t,S_t)-\Delta_tS_t}{B_t}}{the value must be $V_t=f(t,S_t)$}
\dl{V_T=f(T,S_T)=g(S_T)=X}{at $t=T$ nothing is left to average}
\end{deriv}
The $dt$ terms then match too, because $f$ solves the Black--Scholes equation (\cref{card:hedge}). So every claim $g(S_T)$ has a self-financing copy: \textbf{the market is complete}, and $P_t(X)$ is its only arbitrage-free price.
\end{steps}

\paragraph{The $T$-period binomial market.} Bond $B_t=(1+r)^t$; stock $S_t=uS_{t-1}$ or $dS_{t-1}$, $t=1,\dots,T$, with $d<1+r<u$. A path is $\omega=(v_1,\dots,v_T)$ with each $v_t\in\{u,d\}$.
\begin{steps}[resume]
\item First, find every EMM. A measure $Q$ equivalent to $P$ gives every path positive probability, so at each node the up-probability $q$ has $0<q<1$; $q$ may depend on the node. At a node at time $t$:
\begin{deriv}
\dl{\E_Q\big(S_{t+1}\given\F_t\big)=\big(qu+(1-q)d\big)S_t}{$S_{t+1}$ is $uS_t$ or $dS_t$; $S_t$ is known at $t$ (\cref{card:freeze})}
\dl{\E_Q\Big(\frac{S_{t+1}}{(1+r)^{t+1}}\Bigm|\F_t\Big)=\frac{S_t}{(1+r)^t}}{what an EMM needs}
\dl{\iff qu+(1-q)d=1+r}{substitute the first line; cancel $\frac{S_t}{(1+r)^t}$}
\dl{\iff q=\frac{1+r-d}{u-d}=p^*}{one linear equation; $u\ne d$}
\dl{0<p^*<1}{$d<1+r<u$}
\end{deriv}
So there is \textbf{exactly one} EMM $P^*$: up-probability $p^*$ at every node, $P^*(\omega)=(p^*)^j(1-p^*)^{T-j}>0$ for a path with $j$ ups.
\end{steps}

\paragraph{Derivation 3: the binomial market is arbitrage-free.}
\begin{steps}[resume]
\item Any self-financing strategy, discounted. Here $\Delta_{t+1},b_{t+1}$ are chosen at time $t$ (they depend on $S_0,\dots,S_t$) and held over $(t,t+1]$. Write $\tilde V_t=V_t/(1+r)^t$ and $\tilde S_t=S_t/(1+r)^t$:
\begin{deriv}
\dl{\tilde V_t=\Delta_{t+1}\tilde S_t+b_{t+1}}{self-financing: the new holding costs $V_t$; divide by $(1+r)^t$}
\dl{\tilde V_{t+1}=\Delta_{t+1}\tilde S_{t+1}+b_{t+1}}{the same holding one period later}
\dl{\tilde V_{t+1}-\tilde V_t=\Delta_{t+1}\big(\tilde S_{t+1}-\tilde S_t\big)}{subtract: $b_{t+1}$ cancels}
\dl{\E_{P^*}\big(\tilde V_{t+1}-\tilde V_t\given\F_t\big)}{take $\E_{P^*}(\cdot\given\F_t)$ of the last line}
\dl{\quad=\Delta_{t+1}\,\E_{P^*}\big(\tilde S_{t+1}-\tilde S_t\given\F_t\big)}{$\Delta_{t+1}$ is known at $t$ (\cref{card:freeze})}
\dl{\quad=0}{$\tilde S$ is a $P^*$-martingale (Step~10)}
\dl{\E_{P^*}\tilde V_T=\tilde V_0=V_0}{add the $T$ steps (tower property)}
\end{deriv}
\item Suppose there were an arbitrage. Then:
\begin{deriv}
\dl{\E_{P^*}\Big(\frac{V_T}{(1+r)^T}\Big)=V_0=0}{Step~11}
\dl{\E_{P^*}\Big(\frac{V_T}{(1+r)^T}\Big)=\sum_\omega P^*(\omega)\frac{V_T(\omega)}{(1+r)^T}>0}{every $P^*(\omega)>0$, every $V_T(\omega)\ge0$, and one is $>0$}
\end{deriv}
The two lines contradict each other. \textbf{So there is no arbitrage.}
\end{steps}

\paragraph{Derivation 4: the binomial market is complete.}
\begin{steps}[resume]
\item By the uniqueness found in Step~10 and the second key fact, the market is complete. Concretely, copy any claim $X=g(S_T)$ backwards. Let $V_t[s]$ be the value needed at time $t$ when $S_t=s$; start with $V_T[s]=g(s)$. At a node at time $t-1$ with $S_{t-1}=s$, choose $(\Delta_t,b_t)$ to hit both values one period later:
\begin{deriv}
\dl{\Delta_t\,us+b_t(1+r)^t=V_t[us]}{if the stock goes up}
\dl{\Delta_t\,ds+b_t(1+r)^t=V_t[ds]}{if it goes down}
\dl{\Delta_t=\frac{V_t[us]-V_t[ds]}{us-ds}}{subtract the two lines}
\dl{b_t=\frac{uV_t[ds]-dV_t[us]}{(1+r)^t(u-d)}}{put $\Delta_t$ back into either line}
\end{deriv}
\item The cost of this holding at the node:
\begin{deriv}
\dl{V_{t-1}[s]=\Delta_ts+b_t(1+r)^{t-1}}{the holding at time $t-1$ prices}
\dl{\phantom{V_{t-1}[s]}=\frac{p^*V_t[us]+(1-p^*)V_t[ds]}{1+r}}{substitute Step~13 and simplify}
\end{deriv}
\item Put the nodes together. Step back from $t=T$ to $t=1$: at every node the holding chosen there costs exactly the value $V_{t-1}[s]$ that the holding before it is worth at that node. So no money is added or taken out (self-financing), and at $T$ the value is $V_T=g(S_T)=X$. Every claim is attainable: \textbf{the binomial market is complete}, and its price is
\begin{deriv}
\dl{V_0=\E_{P^*}\Big(\frac{g(S_T)}{(1+r)^T}\Big)}{Step~14, repeated $T$ times}
\dl{\phantom{V_0}=\frac{\sum_{j=0}^T\binom Tj(p^*)^j(1-p^*)^{T-j}g\big(u^jd^{T-j}S_0\big)}{(1+r)^T}}{$j$ ups in $\binom Tj$ orders (\cardrow{card:tree}{1})}
\end{deriv}
\end{steps}
\answer{Black--Scholes: $e^{-rt}S_t$ is a martingale only for the drift $r-\frac12\sigma^2$, so the EMM exists and is unique. Binomial: the up-probability must be $p^*=\frac{1+r-d}{u-d}\in(0,1)$ at every node, so again the EMM exists and is unique. In both, existence makes every discounted self-financing value a martingale, so an arbitrage would have mean $0$ and mean $>0$ at once: arbitrage-free. In both, every claim has a self-financing copy ($\Delta_t=\partial_sf$, or the backward one-period copies): complete.}
\end{problem}

\begin{problem}[Q2(d)]
A straddle is a call and a put on the same stock with the same strike and expiry. Plot the payoff of the straddle as a function of the final price $S_T$.
\uses{card:payoff}{4}
\cardsays Since $|x|=x^++x^-$, a call plus a put with the same $K$ pays
\[
(S_T-K)^++(K-S_T)^+=|S_T-K|.
\]
\working
\begin{steps}
\item Below $K$ only the put pays: $K-S_T$, a line of slope $-1$ falling to $0$ at $S_T=K$.
\item Above $K$ only the call pays: $S_T-K$, slope $+1$.
\item Together: a V with its tip at $(K,0)$.
\item Profit: subtract the cost $C+P$. The buyer gains only if $S_T$ ends outside $K\pm(C+P)$: a bet on a big move in either direction.
\end{steps}
\begin{pic}
\begin{tikzpicture}
\begin{axis}[width=8cm,height=5.2cm,xmin=0,xmax=200,ymin=-50,ymax=110,xtick={60,100,140},xticklabels={$K-c$,$K$,$K+c$},ytick={0},xlabel={$S_T$},axis lines=left,ticklabel style={font=\small},legend style={at={(0.5,0.98)},anchor=north,font=\small,draw=none,fill=none}]
  \addplot[exc,very thick] coordinates {(0,100)(100,0)(200,100)};
  \addplot[defc,thick,dashed] coordinates {(0,60)(100,-40)(200,60)};
  \addplot[black!40,thin,forget plot] coordinates {(0,0)(200,0)};
  \legend{payoff $|S_T-K|$,profit $|S_T-K|-c$}
\end{axis}
\end{tikzpicture}
\picnote{Straddle payoff (solid) and profit (dashed), with $c=C+P$ the cost of the call and the put. The profit is positive only outside $K\pm c$.}
\end{pic}
\answer{The payoff is $|S_T-K|$: a V with its tip at $S_T=K$, slope $-1$ to the left and $+1$ to the right.}
\end{problem}

\begin{problem}[Q2(e)]
Explain value at risk (VaR).
\uses{card:normal-prob}{5}
\cardsays The $\alpha$-quantile of a normal $\xi\sim N(\mu,\sigma^2)$ is $q=\mu+z_\alpha\sigma$, where $\Phi(z_\alpha)=\alpha$. For a loss $\xi$, this $q$ is its value at risk at level $\alpha$.
\working
\begin{steps}
\item Definition: fix a horizon (say one day) and a level $\alpha$ (say $99\%$). VaR is the loss level $q$ that is exceeded only with probability $1-\alpha$: $P(\text{loss}\le q)=\alpha$.
\item In words: ``on $99\%$ of days we lose no more than $q$''.
\item Normal loss: $P(\xi\le q)=\Phi\big(\frac{q-\mu}\sigma\big)=\alpha$ gives $\frac{q-\mu}\sigma=z_\alpha$, so $q=\mu+z_\alpha\sigma$, with $z_{0.95}=1.645$ and $z_{0.99}=2.326$.
\item Numbers: loss $\xi\sim N(1,4^2)$ at $\alpha=0.99$: $q=1+2.326(4)=10.31$.
\item What it does not say: how large the loss is on the bad $1\%$ of days.
\end{steps}
\answer{VaR at level $\alpha$ is the $\alpha$-quantile of the loss: the loss exceeded only with probability $1-\alpha$. For a normal loss it is $\mu+z_\alpha\sigma$.}
\end{problem}

%=====================================================================
\section{Question 3: from binomial to Black--Scholes, and a two-period put}

\begin{problem}[Q3(a)]
Using the binomial model, derive the Black--Scholes formula. Then explain the Black--Scholes framework.
\uses{card:tree}{3}
\cardsays In a $T$-period tree, a call costs
\[
C=S_0\sum_{j\ge m}\binom Tj\hat p^{\,j}(1-\hat p)^{T-j}-\frac{K}{(1+r)^T}\sum_{j\ge m}\binom Tj(p^*)^j(1-p^*)^{T-j},
\]
with $p^*=\frac{1+r-d}{u-d}$, $\hat p=\frac{up^*}{1+r}$, and $m$ the least number of ups that ends above $K$.
\uses{card:bsprice}{8}
\cardsays Let the number of steps grow: the tree price tends to the Black--Scholes price.
\working
\begin{steps}
\item Tree price: a path with $j$ ups has weight $(p^*)^j(1-p^*)^{T-j}$ and there are $\binom Tj$ of them, so $C=(1+r)^{-T}\sum_j\binom Tj(p^*)^j(1-p^*)^{T-j}(u^jd^{T-j}S_0-K)^+$. Keeping only $j\ge m$ and splitting off the $S_0$ part (using $up^*+d(1-p^*)=1+r$) gives the box.
\item Shrink the steps: $n$ steps of length $\Delta=\tau/n$, $u=e^{\sigma\sqrt\Delta}$, $d=e^{-\sigma\sqrt\Delta}$, $1+r_n=e^{r\Delta}$.
\item Expand: $p^*=\frac{e^{r\Delta}-e^{-\sigma\sqrt\Delta}}{e^{\sigma\sqrt\Delta}-e^{-\sigma\sqrt\Delta}}\approx\frac12\Big(1+\frac{r-\frac12\sigma^2}{\sigma}\sqrt\Delta\Big)$.
\item This is the walk of \cref{card:scaledwalk} with drift $r-\frac12\sigma^2$, so under $p^*$, $\ln\frac{S_\tau}{S_0}\to N\big((r-\frac12\sigma^2)\tau,\sigma^2\tau\big)$.
\item The $K$ part: $K(1+r_n)^{-n}P^*(S_\tau>K)\to Ke^{-r\tau}\Phi(h-\sigma\sqrt\tau)$.
\item The stock part: $\hat p\approx\frac12\big(1+\frac{r+\frac12\sigma^2}{\sigma}\sqrt\Delta\big)$, a walk with drift $r+\frac12\sigma^2$, so the first sum tends to $\Phi(h)$, with $h=\frac{\ln(S_0/K)+(r+\frac12\sigma^2)\tau}{\sigma\sqrt\tau}$.
\item Together: $C=S_0\Phi(h)-Ke^{-r\tau}\Phi(h-\sigma\sqrt\tau)$.
\end{steps}
\textbf{The framework} has these parts.
\begin{itemize}
\item A bond $B_t=e^{rt}$ and a stock $S_t=S_0e^{at+\sigma W_t}$ (geometric Brownian motion), with $r$ and $\sigma>0$ constant.
\item Continuous trading, no costs, no dividends, short selling allowed.
\item The risk-neutral measure exists and is unique (drift $r-\frac12\sigma^2$), so the market is arbitrage-free and complete, and every claim costs $e^{-r\tau}\E_{P^*}[g(S_T)\given\F_t]$.
\end{itemize}
\answer{The binomial price tends to $C=S_0\Phi(h)-Ke^{-r\tau}\Phi(h-\sigma\sqrt\tau)$, the Black--Scholes formula.}
\end{problem}

\begin{problem}[Q3(b)]
Two-period binomial market, $u=\frac53$, $d=\frac13$, $r=0$, $S_0=60$. Price a put with maturity $T=2$ and strike $K=80$.
\uses{card:tree}{1}
\cardsays Start at the end with the payoffs, then step back: each node is the risk-neutral average of the next two values, divided by $1+r$,
\[
V_t[s]=\frac{p^*\,V_{t+1}[us]+(1-p^*)\,V_{t+1}[ds]}{1+r},\qquad V_T[s]=g(s).
\]
\working
\begin{steps}
\item $p^*=\frac{1+r-d}{u-d}=\frac{1-\frac13}{\frac53-\frac13}=\frac{2/3}{4/3}=\frac12$. Guard: $\frac13<1<\frac53$ holds.
\item Stock tree. Time 1: $100$ and $20$. Time 2: $\frac{25}9(60)=166.67$, $\frac59(60)=33.33$, $\frac19(60)=6.67$.
\item Put payoffs $(80-S_2)^+$: $0$, $46.67$, $73.33$.
\item Node $100$: $V_1=\frac12(0)+\frac12(46.67)=23.33$.
\item Node $20$: $V_1=\frac12(46.67)+\frac12(73.33)=60$.
\item Today: $V_0=\frac12(23.33)+\frac12(60)=41.67$.
\end{steps}
\answer{The put costs $\frac{125}3\approx41.67$.}
\end{problem}

%=====================================================================
\section{Question 4: a two-period call, its hedge, and the put}

Market: $T=2$, $u=1.35$, $d=0.6$, $r=0.15$, $S_0=400$, strike $K=500$.

\begin{problem}[Q4(i)]
Price a call with maturity $T=2$ and strike $K=500$.
\uses{card:tree}{1}
\cardsays Backward induction, $V_t[s]=\frac{p^*V_{t+1}[us]+(1-p^*)V_{t+1}[ds]}{1+r}$, starting from the payoffs.
\working
\begin{steps}
\item $p^*=\frac{1.15-0.6}{1.35-0.6}=\frac{0.55}{0.75}=\frac{11}{15}=0.7333$, so $1-p^*=0.2667$. Guard: $0.6<1.15<1.35$ holds.
\item Stock tree. Time 1: $540$ and $240$. Time 2: $729$, $324$, $144$.
\item Call payoffs $(S_2-500)^+$: $229$, $0$, $0$.
\item Node $540$: $V_1=\frac{0.7333(229)+0.2667(0)}{1.15}=\frac{167.93}{1.15}=146.03$.
\item Node $240$: $V_1=\frac{0+0}{1.15}=0$.
\item Today: $V_0=\frac{0.7333(146.03)+0}{1.15}=\frac{107.09}{1.15}=93.12$.
\end{steps}
\answer{The call costs $93.12$.}
\end{problem}

\begin{problem}[Q4(ii)]
Construct a replicating portfolio for this option.
\uses{card:tree}{5}
\cardsays At each node hold $\Delta$ shares and the rest of the value in the bond:
\[
\Delta=\frac{V_{t+1}[us]-V_{t+1}[ds]}{us-ds},\qquad\text{bond money}=V_t[s]-\Delta s.
\]
\working
\begin{steps}
\item Today ($s=400$): $\Delta_1=\frac{146.03-0}{540-240}=\frac{146.03}{300}=0.4868$; bond $=93.12-0.4868(400)=93.12-194.71=-101.59$ (borrowed).
\item Node $540$: $\Delta_2=\frac{229-0}{729-324}=\frac{229}{405}=0.5654$; bond $=146.03-0.5654(540)=146.03-305.33=-159.30$.
\item Node $240$: $\Delta_2=\frac{0-0}{324-144}=0$; bond $=0-0=0$.
\end{steps}
\answer{Today: $0.4868$ shares and borrow $101.59$. After an up-move: $0.5654$ shares and borrow $159.30$. After a down-move: hold nothing.}
\end{problem}

\begin{problem}[Q4(iii)]
Directly verify that this strategy is self-financing.
\uses{card:tree}{6}
\cardsays Self-financing: at each rebalancing, the old holding valued at the new prices equals the cost of the new holding.
\working
\begin{steps}
\item The borrowing grows to $-101.59(1.15)=-116.82$ at time 1.
\item Up ($S_1=540$): old holding $0.4868(540)-116.82=262.85-116.82=146.03$. New holding costs $0.5654(540)-159.30=146.03$. Equal.
\item Down ($S_1=240$): old holding $0.4868(240)-116.82=116.82-116.82=0$. New holding costs $0$. Equal.
\item Time 2 from node $540$: the borrowing grows to $-159.30(1.15)=-183.20$. At $729$: $0.5654(729)-183.20=412.20-183.20=229$. At $324$: $0.5654(324)-183.20=183.20-183.20=0$. Exactly the call.
\end{steps}
\answer{At each node the old holding is worth exactly what the new one costs, so no money goes in or out: self-financing.}
\end{problem}

\begin{problem}[Q4(iv)]
Price a put with the same maturity and strike.
\uses{card:parity}{1}
\cardsays Put--call parity solved for the put: $P=C-S_0+KD_T$, with $D_T=(1+r)^{-T}$.
\working
\begin{steps}
\item Discount: $D_2=1.15^{-2}=\frac1{1.3225}$, so $KD_2=\frac{500}{1.3225}=378.07$.
\item Substitute: $P=93.12-400+378.07=71.19$.
\end{steps}
\answer{The put costs $71.19$.}
\end{problem}

%=====================================================================
\section{Question 5: deriving the Greeks}

\begin{problem}[Q5]
Derive the sensitivities delta, gamma, theta, vega and rho of a European call in the Black--Scholes framework.
\uses{card:bsprice}{6}
\cardsays With $\tau=T-t$ and $k=h-\sigma\sqrt\tau$, the call is $C=S\Phi(h)-Ke^{-r\tau}\Phi(k)$, and
\[
\begin{gathered}
\Delta=\Phi(h),\qquad\Gamma=\frac{\varphi(h)}{S\sigma\sqrt\tau},\qquad\text{vega}=S\varphi(h)\sqrt\tau,\\
\rho=K\tau e^{-r\tau}\Phi(k),\qquad\theta=-\frac{S\varphi(h)\sigma}{2\sqrt\tau}-rKe^{-r\tau}\Phi(k).
\end{gathered}
\]
The trick is one identity that makes the messy terms cancel.
\working
\begin{steps}
\item The identity $S\varphi(h)=Ke^{-r\tau}\varphi(k)$. Since $\varphi(x)=\frac{e^{-x^2/2}}{\sqrt{2\pi}}$,
\[
\begin{aligned}
\frac{\varphi(k)}{\varphi(h)}&=e^{(h^2-k^2)/2}=e^{(h-k)(h+k)/2}=e^{\sigma\sqrt\tau(2h-\sigma\sqrt\tau)/2}\\
&=e^{\sigma\sqrt\tau\,h-\frac12\sigma^2\tau}=e^{\ln(S/K)+(r+\frac12\sigma^2)\tau-\frac12\sigma^2\tau}=\frac SKe^{r\tau}.
\end{aligned}
\]
Multiply both sides by $Ke^{-r\tau}\varphi(h)$.
\item How $k$ moves with $h$: $k=h-\sigma\sqrt\tau$, so $\partial k=\partial h$ for $S$ and $r$; $\partial_\sigma k=\partial_\sigma h-\sqrt\tau$; $\partial_\tau k=\partial_\tau h-\frac{\sigma}{2\sqrt\tau}$.
\item \textbf{Delta.} $\partial_SC=\Phi(h)+\big[S\varphi(h)-Ke^{-r\tau}\varphi(k)\big]\partial_Sh=\Phi(h)$, because the bracket is $0$ by Step 1.
\item \textbf{Gamma.} $\partial_S\Phi(h)=\varphi(h)\,\partial_Sh$, and $\partial_Sh=\frac{1}{S\sigma\sqrt\tau}$ (only $\ln S$ contains $S$). So $\Gamma=\frac{\varphi(h)}{S\sigma\sqrt\tau}$.
\item \textbf{Vega.} $\partial_\sigma C=S\varphi(h)\partial_\sigma h-Ke^{-r\tau}\varphi(k)\big(\partial_\sigma h-\sqrt\tau\big)$. The $\partial_\sigma h$ terms cancel by Step 1, leaving $Ke^{-r\tau}\varphi(k)\sqrt\tau=S\varphi(h)\sqrt\tau$.
\item \textbf{Rho.} $\partial_rC=\big[S\varphi(h)-Ke^{-r\tau}\varphi(k)\big]\partial_rh+\tau Ke^{-r\tau}\Phi(k)=K\tau e^{-r\tau}\Phi(k)$.
\item \textbf{Theta.} First in $\tau$: $\partial_\tau C=S\varphi(h)\partial_\tau h-Ke^{-r\tau}\varphi(k)\big(\partial_\tau h-\frac\sigma{2\sqrt\tau}\big)+rKe^{-r\tau}\Phi(k)=\frac{S\varphi(h)\sigma}{2\sqrt\tau}+rKe^{-r\tau}\Phi(k)$. Time $t$ runs the other way ($\tau=T-t$), so $\theta=\partial_tC=-\partial_\tau C$.
\item Numbers ($S=80$, $K=85$, $r=0.05$, $\sigma=0.15$, $\tau=1$, from Q1(e)): $\Delta=0.5017$, $\Gamma=0.0332$, vega $=80(0.3989)=31.92$, $\rho=85e^{-0.05}(0.4420)=35.74$, $\theta=-2.39-1.79=-4.18$.
\end{steps}
\answer{$\Delta=\Phi(h)$, $\Gamma=\frac{\varphi(h)}{S\sigma\sqrt\tau}$, vega $=S\varphi(h)\sqrt\tau$, $\rho=K\tau e^{-r\tau}\Phi(h-\sigma\sqrt\tau)$, $\theta=-\frac{S\varphi(h)\sigma}{2\sqrt\tau}-rKe^{-r\tau}\Phi(h-\sigma\sqrt\tau)$.}
\end{problem}
